\documentclass[../../main.tex]{subfiles}
\begin{document}

\begin{flushright}
  \footnotesize\textit{【自動文字起こし・要確認】}
\end{flushright}

% Source: handwritten.pdf 第5問．
{\bf [解答]}

\paragraph{(1)}
$a_n=3^{2^n}-1$とおく．$a_n$が$2^{n+2}$で割り切れ，$2^{n+3}$で割り切れないことを帰納法で示す．
$n=1$のとき$a_1=8=2^3$だから成立する．$n=k$で成立すると仮定する．
\begin{align}
  a_{k+1}=(3^{2^k}+1)(3^{2^k}-1).
\end{align}
ここで
\begin{align}
  3^{2^k}+1&\equiv 1^{2^k}+1\equiv0\pmod{2},\\
  3^{2^k}+1&\equiv(-1)^{2^k}+1\equiv2\pmod{4}.
\end{align}
よって第1因子は2で割り切れ，4で割り切れない．仮定より$a_{k+1}$は$2^{k+3}$で割り切れ，$2^{k+4}$で割り切れない．帰納法により示された．

\paragraph{(2)}
$m=A2^n$（$A$は奇数，$n\geqq1$）と表す．
\begin{align}
  3^m-1&=(3^{2^n})^A-1=(a_n+1)^A-1\\
  &=a_n\left\{a_n^{A-1}+\binom{A}{1}a_n^{A-2}+\cdots+\binom{A}{A-2}a_n+A\right\}.
\end{align}
中括弧内は奇数だから，(1)より$3^m-1$は$2^{n+2}$で割り切れ，$2^{n+3}$で割り切れない．したがって必要な条件は
\begin{align}
  A2^n\leqq n+2.
\end{align}
$n\geqq3$では$A2^n\geqq n+3$である．実際，$n=3$では$A2^3\geqq8\geqq6$であり，$A2^k\geqq k+3$なら
\begin{align}
  A2^{k+1}\geqq2(k+3)\geqq k+4
\end{align}
だから帰納法で従う．よって$n\geqq3$は不適．
$n=1$では$2A\leqq3$より$A=1$，$m=2$．$n=2$では$4A\leqq4$より$A=1$，$m=4$．以上より
\begin{align}
  m=2,4.
\end{align}

\end{document}
